\documentclass[11pt]{article}

\usepackage[margin=1in]{geometry}
\usepackage{booktabs}
\usepackage{tabularx}
\usepackage{array}
\usepackage{amsmath}
\usepackage{xeCJK}
\usepackage{microtype}
\usepackage{hyperref}
\usepackage{enumitem}

\setCJKmainfont{SimSun}
\setmainfont{Times New Roman}
\setlist{nosep}
\newcolumntype{Y}{>{\raggedright\arraybackslash}X}

\title{SRTP 原理方向证据合成报告}
\author{SRTP Optical Mechanism Analysis}
\date{2026-06-22}

\begin{document}
\maketitle

\section{当前论文主线建议}

本项目当前最稳妥的投稿主线应收束为：面向反光工业微结构表面的焦栈三维重建中，传统 DFF 的失效主要来自焦度响应不确定、局部饱和/高亮和弱信号区域，而不仅是简单的采集漂移或单一眩光几何风险。本文通过 Simulation-to-Real 数据构造、DFF/GADFF 先验、focus-confidence 诊断和置信度加权训练策略，构建一个先验引导的稳健重建框架。

这条主线的关键优点是把 ``Simulation-to-Real'' 定位为解决真实标注不足和退化覆盖不足的手段，同时把论文的科学问题放在“反光表面焦栈中 DFF 为什么失效、如何被识别和缓解”上。

\section{证据等级总表}

\small
\begin{tabularx}{\linewidth}{p{0.07\linewidth}Y Y p{0.12\linewidth}Y}
\toprule
编号 & 可支撑的机制判断 & 主要证据 & 证据等级 & 主要限制 \\
\midrule
E1 & DFF 失败更稳定地对应 top-2 focus margin 低，而非几何 glare risk 单独决定 & 135 个合成条件；\texttt{low\_margin} mean AUC 0.6412，std 0.0133，AUC>0.55 rate 1.00；\texttt{risk\_mean/risk\_max} mean AUC 约 0.51 & 强机制证据 & 基于合成数据，需避免写成真实 MAE 结论 \\
E2 & 真实焦栈中的不可靠区域具有多种焦度曲线形态，不宜合并成单一失败类别 & 7 组真实焦栈；清晰单峰 49.5\%，平坦歧义 25.5\%，暗弱信号 15.2\%，多峰 4.2\%，局部跳变 3.8\%，高亮饱和 1.8\% & 中强真实诊断证据 & 无真实高度 GT，属于 no-reference morphology \\
E3 & 100um 真实焦栈的局部 DFF 跳变难以由小幅全局层间平移单独解释 & 最大层间平移 0.1118 px，中位数 0；配准后 peak layer 变化 4.68\%；spike proxy Pearson 0.9478；quality proxy Pearson 0.9825 & 中强排除性证据 & 只排除全局平移，不能排除非刚性、倍率、照明角变化 \\
E4 & 置信度加权伪标签训练在多退化混合场景中更有价值 & 96 个受控合成条件；总体 MAE 0.0883 到 0.0857，胜率 71.9\%；mixed 模式改善 7.2\%，胜率 100.0\% & 中等训练策略代理证据 & 轻量学生模型，不等价于完整 FocusResUNet \\
E5 & DFF/GADFF 先验仍是当前模型体系中最稳定的几何基础 & 已有 ablation：移除 DFF/GADFF prior 后 MAE 133.4808，高于 Full S2R-FocusNet 109.2209；但 w/o focal difference 和 w/o glare cue 更优 & 中等模型消融证据 & 完整模型未占优，需要写成“DFF prior 有价值，融合设计需重构” \\
\bottomrule
\end{tabularx}
\normalsize

\section{可直接写进论文的主张}

\subsection{Focus-margin confidence 是当前最可靠的 DFF failure 指示器}

可写主张：In the controlled synthetic sweep, the top-2 focus margin provides the most stable no-reference indicator of DFF failure regions. The low-margin score achieves a mean AUC of 0.6412 with a standard deviation of 0.0133 over 135 conditions and remains above 0.55 in all tested cases.

边界：不应写成 \texttt{glare risk} 无用；它仍可解释反光形成条件。也不应写成 \texttt{low\_margin} 已在真实高度 GT 上证明有效。

\subsection{真实焦栈不可靠性是多形态问题}

可写主张：Real focus stacks exhibit multiple focus-response morphologies rather than a single failure mode. Across seven real stacks, only 49.5\% of pixels are classified as confident single-peak responses, while the rest are distributed among flat ambiguous responses, dark low-signal regions, multi-peak competition, local peak-layer spikes, and saturated highlights.

边界：这些类别不能写成真实错误类别，比例也不等同于真实深度失败比例。

\subsection{全局层间平移不是当前 100um 焦栈不稳定性的主要解释}

可写主张：For the 100um key-texture real stack, the estimated global inter-layer shift is below 0.12 pixels. After translation compensation, only 4.68\% of pixels change their selected DFF peak layer, and the low-margin AUC for identifying spike-proxy pixels remains nearly unchanged.

边界：该证据不能排除非刚性形变、倍率变化、入瞳/照明变化或离焦导致的结构外观变化。

\subsection{Confidence-weighted pseudo-labeling 是有条件收益的训练策略}

可写主张：In a controlled pseudo-labeling study, confidence-weighted supervision reduces the average MAE of a lightweight student from 0.0883 to 0.0857 over 96 conditions. The benefit is concentrated under mixed degradation, where the relative improvement reaches 7.2\% with a 100.0\% win rate.

边界：该结果不能替代完整 FocusResUNet 训练验证，也不能声称 weak-texture 场景有明显收益。

\section{建议的论文贡献重写}

\begin{enumerate}
  \item \textbf{反光表面 DFF 失效的焦度响应机制分析。} We analyze DFF failure in reflective focus-stack reconstruction from the perspective of focus-response ambiguity, showing that low top-2 focus margin is a more stable failure indicator than geometric glare risk alone.
  \item \textbf{面向真实退化的 Simulation-to-Real 数据与诊断先验。} We design a simulation-to-real data construction strategy that covers not only surface height variation, but also focus-curve morphologies observed in real stacks, including flat responses, multi-peak competition, saturated highlights, and low-signal regions.
  \item \textbf{置信度感知的先验引导重建训练策略。} We formulate DFF/GADFF priors as confidence-aware observations and use focus-confidence maps to guide pseudo-label weighting and prior fusion, instead of treating all DFF-derived targets as equally reliable.
\end{enumerate}

\section{当前最安全的方法公式}

建议将 quality prior 写成分层形式：

\[
Q_{\mathrm{margin}}(x)=1-M_{\mathrm{focus}}(x),\quad
R_{\mathrm{glare}}(x)=\mathrm{geometry\text{-}based\ glare\ risk},\quad
P_{\mathrm{sat}}(x)=\mathrm{saturation\ persistence}.
\]

其中，$Q_{\mathrm{margin}}$ 是训练权重和 DFF 失败诊断的主信号；$R_{\mathrm{glare}}$ 是物理解释、区域分层和仿真条件控制信号；$P_{\mathrm{sat}}$ 用于区分饱和/高亮区域和普通低置信度区域。

对应训练表述为：

\[
L=\mathrm{mean}\left(w(x)\cdot |H_{\mathrm{pred}}(x)-H_{\mathrm{target}}(x)|\right),\quad
w(x)=f(Q_{\mathrm{margin}}(x),P_{\mathrm{sat}}(x),R_{\mathrm{glare}}(x)).
\]

论文中应强调 $w(x)$ 是 confidence-aware weighting，而不是把 $R_{\mathrm{glare}}$ 直接等价为 error mask。

\section{需要继续补强的缺口}

\small
\begin{tabularx}{\linewidth}{Y Y Y}
\toprule
缺口 & 为什么重要 & 可执行下一步 \\
\midrule
完整模型中的 confidence-weighted loss 尚未训练验证 & 代理实验不能替代神经网络主实验 & 在现有 ablation runner 中增加 margin-weighted loss 变体，做 matched split smoke + full candidate \\
w/o focal difference 和 w/o glare cue 反而优于 full model & 说明现有融合方式可能稀释稳定先验 & 将 glare/focal cue 从直接拼接改为 gating 或 confidence head \\
真实样本缺少高度 GT & 限制真实结果只能写 no-reference diagnosis & 选择少量 ROI 做人工标注、重复采集或外部轮廓仪/共聚焦验证 \\
真实焦栈只做了全局平移敏感性 & 仍可能存在非刚性、倍率、照明路径变化 & 增加局部块配准或 feature-flow 诊断，但只作为 sanity check \\
SOTA 对比仍需更新 & 投稿时必须避免过老 baseline & 保留 Depth Anything/Depth Anything V2 等 modern prior 作为相关工作，不一定作为直接 baseline \\
\bottomrule
\end{tabularx}
\normalsize

\section{下一步实验建议}

优先级最高的下一步是在 \texttt{tmp/ablation\_results} 内新增一个 \texttt{confidence\_weighted\_loss} 训练变体。保持现有 matched split 和 protected runner，不改动主项目交付代码。先做 smoke 验证，再做 4-epoch matched full candidate。指标重点看 global MAE、high-risk MAE、spike-region MAE 和 p90 error。若结果优于 full model，再进入 seed repeat；若结果不优于 full model，则把当前证据写成 training motivation，而不是主方法收益。

\section{Paper-Ready Paragraph}

Our diagnostic analysis suggests that the degradation of DFF in reflective focus stacks is governed more directly by focus-response ambiguity than by glare-prone geometry alone. In a 135-condition synthetic sweep, the low top-2 focus margin achieves a mean AUC of 0.6412 for identifying top DFF failure regions and remains above 0.55 in all conditions, whereas geometry-based glare-risk scores are close to random ranking. Real focus stacks further show that unreliable DFF observations are heterogeneous: across seven stacks, only 49.5\% of pixels exhibit confident single-peak responses, while the remaining regions are distributed among flat ambiguous responses, dark low-signal areas, multi-peak competition, local peak-layer spikes, and saturated highlights. A registration sensitivity test on the 100um real stack also shows that the observed local DFF layer-selection jumps cannot be explained by small global inter-layer translation alone. These findings motivate treating DFF/GADFF priors as confidence-aware observations and using focus-confidence maps for loss weighting, pseudo-label selection, and region-wise analysis.

\end{document}
